% AUTHORED
\section{Multi-layer perceptron}\label{sec:04}
\conceptbadge{04}

A multi-layer perceptron (MLP) is a stack of dense layers separated by nonlinear
activation functions. A two-layer version can be written as
\[
  H = \phi(XW_1 + b_1),
  \qquad
  Y = HW_2 + b_2,
\]
where $\phi$ might be ReLU, tanh, or GELU. The hidden matrix $H$ is a learned
representation of the input.

\subsection*{The role of the activation}
Concept 02 showed that linear layers alone collapse into one linear layer. The
activation prevents that collapse. ReLU is the simplest common example:
\[
  \operatorname{ReLU}(x) = \max(0,x).
\]
It keeps positive values and zeroes negative values. This piecewise-linear bend
is enough to let a network form decision regions that a single linear unit
cannot.

\subsection*{Batched shapes}
For a batch $X \in \mathbb{R}^{B \times d_{\text{in}}}$ and hidden width
$d_{\text{hid}}$, the first layer parameters have shapes
\[
  W_1 \in \mathbb{R}^{d_{\text{in}} \times d_{\text{hid}}},
  \qquad
  b_1 \in \mathbb{R}^{d_{\text{hid}}}.
\]
The hidden activations have shape $B \times d_{\text{hid}}$. If the output
dimension is $d_{\text{out}}$, then
$W_2 \in \mathbb{R}^{d_{\text{hid}} \times d_{\text{out}}}$ and
$Y \in \mathbb{R}^{B \times d_{\text{out}}}$. Writing these shapes before
coding avoids most dense-layer bugs.

\subsection*{Features, not just labels}
Hidden units learn features useful for the loss. In a classifier those features
help separate classes. In a language model they become token-level
representations that support next-token prediction. The MLP inside a transformer
does not mix tokens; it transforms each token's channel vector independently,
which complements attention's token-mixing role.

\subsection*{Training}
The forward pass produces logits, the loss measures the error, and backprop
computes gradients for every weight and bias. In P2 you will turn the pieces
into reusable primitives: linear layers, activations, losses, and optimizers.
The module exercise isolates just the forward computation of an affine layer
plus ReLU so the shape mechanics are solid before gradients are added.

\medskip
\noindent\textit{Go deeper:} Goodfellow, Bengio \& Courville, \emph{Deep
Learning}, Ch.~6 (\texttt{https://www.deeplearningbook.org}) develops
feed-forward networks in depth; Karpathy's \texttt{makemore}
(\texttt{https://github.com/karpathy/makemore}) builds an MLP language model
step by step.
